\scanpage{146}
% archive foliation: 145

évident, d'ailleurs pleinement fidèle
\[
\underline h^{\unsure{\sharp|*|\natural}}:\mathcal V\longrightarrow\mathcal M^{\unsure{\sharp|*|\natural}}
\]
en posant \(\underline h^{\unsure{\sharp|*|\natural}}(X)=(X,\operatorname{id}_X)\).
Pour toute catégorie
abélienne \(\mathcal C\), le foncteur \(\varphi:T\longmapsto T\circ h\)
\[
\underline{\operatorname{Fonc\,add}}(\mathcal M,\mathcal C)
\xrightarrow{\varphi}
\underline{\operatorname{Fonc\,add}}(\mathcal V,\mathcal C)
\]
est une \deleted{isomorphisme} équivalence. Pour le voir, on
se construit un foncteur quasi-inverse
\deleted{(1) pleinement fidèle. Un}
\[
\psi:\underline{\operatorname{Fonc\,add}}(\mathcal V,\mathcal C)
\longrightarrow
\underline{\operatorname{Fonc\,add}}(\mathcal M,\mathcal C)
\]
en associant à \unsure{un} foncteur \inserted{additif}
\(H:\mathcal V\mathbin{-}\mathcal C\)
son extension canonique
\[
\psi(H)(X,\pi)=\operatorname{Im}H(\pi)
\]
avec la loi fonctorielle évidente. On voit
aussitôt que l'on a des isom can
\[
\begin{aligned}
\varphi(\psi\unsure{)}(H)&\simeq H\\
\psi\varphi(T)&\simeq T
\end{aligned}
\]
Pour l'instant, cette construction et propriété
universelle de \(\mathcal M\) n'utilise pas les hyp
précédentes faites sur \(\mathcal V\); elle répond au Pb
universel de \uline{rendre tous les projecteurs
scindables} \inserted{(d'une catégorie additive donnée)}.
Mais les hypothèses faites
sur \(\mathcal V\) vont impliquer que \(\mathcal M\) est en
fait une catégorie \inserted{\uline{abélienne}} \uline{semi-simple}
(i.e. tout objet de

\margin{il suffit que \(\mathcal C\) soit une catégorie additive où
les projecteurs sont scindables}
